\scanpage{18}
\mspage{\deleted{\unsure{58}}4}

l'exposé \deleted{ultérieur}\inserted{$\overline{\underline{\mathrm{XIV}}}$} \deleted{que la réciproque est vraie}\inserted{que la réciproque est vraie} lorsque $S$ est \deleted{\unsure{noethérien}}\inserted{artinien}, \deleted{\unsure{dimension quelconque}}\deleted{\unsure{artinien}} \inserted{ou} \deleted{\unsure{et}} lorsque $S$ est un schéma local et $G$ est ``réductif'' : dans ce cas, tout tore de $G$ est contenu dans un tore maximal.

\uline{Définition} 1.5. Soit $G$ un groupe algébrique sur un corps $k$. \deleted{\unsure{suivant}}
On appelle rang réductif (resp. rang nilpotent, resp. rang unipotent \unsure{\ldots}) de $G$ le rang réductif (resp. \ldots) de $G_{\bar k}$, où $\bar k$ est une clôture algébrique de $k$.

On notera qu'en vertu de 1.2. et de la commutation de la formation de $\underline{\operatorname{Centr}}_G(T)$ avec l'extension de la base, \deleted{les} les \deleted{rangs} notions de rang introduites dans 1.5. sont \uline{invariantes par extension du corps de base}; d'autre part, pour $k$ algébriquement clos, elles coïncident avec celles introduites au début du présent numéro.

\uline{Remarque} 1.6. Il n'est pas difficile de construire un schéma en groupes affine et lisse $G$ sur le spectre $S$ d'un anneau de valuation discrète, dont la fibre générale soit isomorphe à $\underline G_m$, et la fibre spéciale isomorphe à $\underline G_a$. Un tel $G$ ne contient aucun tore sauf le tore trivial \inserted{$T$} (réduit au sous-groupe unité), qui n'est évidemment pas tore maximal. De façon précise, dans la fibre spéciale $G_0=\underline G_{a\,k}$, $T_0$ est bien tore maximal, mais dans la fibre générale $G_1=\underline G_{m\,K}$, $T_1$ n'est plus tore maximal (NB $k$ est le corps résiduel, $K$ le corps des fractions de l'anneau de valuation envisagé). On voit aussi sur cet exemple que le rang réductif de $G_s$ ($s\in S$) n'est pas fonction continue de $s$. On a cependant les résultats suivants:

\uline{Théorème} 1.7. Soient $S$ un préschéma, $G$ un $S$-préschéma en groupes affine et lisse sur $S$. Pour tout $s\in S$, considérons $\rho_r(s)=\rho_r(G_s)$ et $\rho_n(s)=\rho_n(G_s)$, les rangs réductif\deleted{\unsure{s}} et nilpotent de $G_s$ (1.5.). Avec ces notations, on a ce qui suit:

\mspage{17}
